\chapter{Discussion and Limitations}

Chapter 8 presented the evaluation evidence. This chapter explains how that
evidence should be read and where the design tradeoffs are. The central
interpretation is clear: Raft can remain in Paxos's performance class inside
Mako when the integration preserves Mako's execution, replication, and replay
pipeline.

\section{What the Results Mean}

The successful Raft configurations should be interpreted as pipeline results,
not protocol-only results. Raft supplies leader election and replicated log
agreement, but Mako's throughput also depends on busy leader-side workers,
enough replication parallelism, correct delivery to follower replay paths, and
followers that replay close to the committed history.

Within the evaluated in-memory runs, Multi-Raft stays close to Paxos while
preserving the original partition-level replication shape, and Single-Raft
with ReplayPool keeps follower replay close to leader-side committed
throughput. In those configurations, replication and follower replay are not
the dominant bottleneck preventing Mako's worker threads from doing useful
work. The result gives a strong systems answer for the evaluated workload,
machine, thread range, and storage models.

This distinction matters because the thesis evaluates a complete implementation
path inside Mako, not only a consensus algorithm in isolation. When the
surrounding systems choices are right, Raft does not prevent the application
layer from scaling in the tested range. When those choices are wrong, as in
Single-Raft without enough replay parallelism, a correct consensus protocol can
still produce an unhealthy replicated system.

\section{Multi-Raft Versus Single-Raft}

Multi-Raft is the conservative design. It keeps Mako close to the original
Paxos-shaped topology by giving separate worker-lane streams separate Raft
groups. This makes the comparison to Paxos straightforward, but it preserves
more replication structure inside the process.

Single-Raft is the consolidation design. It places all worker-lane commands
into one consolidated Raft log, with one local replica of that log in each
process. This simplifies the replication topology, but it concentrates work
after commit. The tradeoff is therefore not simply ``many Raft groups'' versus
``one Raft group.'' It is whether the full pipeline preserves enough
parallelism after the consensus topology is changed.

The designs emphasize different strengths. Multi-Raft is the strongest choice
when the priority is preserving Mako's original performance shape. Single-Raft
is the attractive consolidation choice when reducing replication structure
matters, and it becomes credible once ReplayPool keeps follower replay moving.
That is why the evaluation treats follower replay progress as a first-class
metric rather than as an implementation detail. The replay-lag proof in
Chapter 8 adds a stronger end-of-run check: in the measured benchmark, 1:1
ReplayPool drains the records counted by the leader-side accounting path,
while the no-ReplayPool design leaves a large unreplayed gap at high worker
counts.

\section{Preferred Leadership as a Placement Mechanism}

Preferred leader election is a Mako-specific placement mechanism around Raft.
It aligns Raft leadership with the leader-side role where Mako submits
transaction log records, while preserving Raft's ordinary voting, log matching,
and commit rules.

This requirement applies to both Raft topologies. Multi-Raft needs preferred
leadership so each Raft group normally accepts transaction records from the
expected process. Single-Raft needs the same placement property for the
consolidated group, and then additionally needs ReplayPool because
consolidation creates replay pressure after commit.

The important safety boundary is explicit. Preferred leadership chooses among
safe Raft leaders during normal operation; voting, log freshness, and commit
rules remain the authority. If the preferred replica is unavailable, backup
replicas can still lead. If the preferred replica returns, leadership can move
back only after catch-up.

\section{ReplayPool and the Application-Level Ordering Caveat}

ReplayPool changes where follower-side database replay runs after Raft has
already committed a record. The key design boundary is wall-clock replay
completion order: a later committed Raft entry for one worker lane may finish
replay before an earlier committed entry for another lane if the two entries
are assigned to different replay workers. This is a deliberate Mako-specific
integration choice, not a generic rule for every replicated state machine.

ReplayPool is appropriate for this Mako integration under the ordering
assumptions evaluated here. Records carry the lane information needed for
routing. Same-lane records stay on the same replay path, so a lane's own order
is preserved. Mako's watermark checks and atomic compare-and-swap style state
updates prevent unsafe out-of-order state changes. The boundary is clear: Raft
provides commit order, while Mako is responsible for making parallel
application replay safe for Mako's data model.

\section{Persistence and Simulated-Disk Scope}

The simulated disk gives controlled, repeatable storage pressure and records
the modeled bytes and writes charged to each backend. That accounting is why
the Cloud-SSD-like flattening can be interpreted with storage-work evidence
rather than only with a throughput curve.

The NVMe-like and Cloud-SSD-like settings are controlled models used to compare
backends under the same assumptions. This keeps the result focused: the thesis
evaluates how each backend behaves under the same modeled storage pressure,
while full device, filesystem, cloud-volume, and tail-latency behavior remain
part of a production-storage evaluation.

Within that scope, the simulated-disk counters are central evidence. They turn
a qualitative curve shape into a measurable explanation: the evaluation shows
how many bytes and writes were charged while the throughput changed. This is
why the disk results support a controlled storage-model claim.

\section{Scope of the Result}

The thesis claim is intentionally concrete: in the evaluated Mako setting, Raft
can replace the original Paxos backend without leaving Paxos's performance
class when topology, leader placement, and replay parallelism preserve Mako's
pipeline. Multi-Raft and Single-Raft answer that question in different ways;
the thesis does not need one design to dominate the other in every deployment.

The result is tied to the measured workload, machine, thread range, and
benchmark harness. It is strong evidence for the thesis claim in this setting
and a clear guide for how to integrate Raft into a high-throughput transaction
pipeline. Broader workloads, physical storage stacks, and failure-recovery
performance under active load are natural next evaluations. The replay-lag
proof establishes end-of-run drain behavior for this benchmark; per-entry
latency bounds across arbitrary workloads are a separate measurement question.
